% Figure for Oscillations, Waves and Optics §A4.1 (a sinusoidal light wave entering and leaving a slab of index n = 1.5: same frequency, wavelength shortened to lambda_0/n inside; reflection neglected).
\begin{tikzpicture}[font=\small, >={Stealth[length=2.0mm]}]
  \begin{axis}[nfig, width=12.5cm, height=4.2cm, xmin=0, xmax=9.6, ymin=-1.6, ymax=1.9,
      axis lines=middle, xtick=\empty, ytick=\empty, xlabel={$x$}, ylabel={$E$},
      xlabel style={anchor=west}, ylabel style={anchor=south}, clip=false]
    \fill[black!8] (axis cs:3,-1.35) rectangle (axis cs:6,1.35);
    \addplot[figB, domain=0:3, samples=150] {sin(deg(2*pi*x))};
    \addplot[figR, domain=3:6, samples=200] {sin(deg(2*pi*(3+1.5*(x-3))))};
    \addplot[figB, domain=6:9.4, samples=150] {sin(deg(2*pi*(7.5+(x-6))))};
    \node[black!60, font=\scriptsize] at (axis cs:1.5,-1.5) {vacuum (air), $n = 1$};
    \node[black!60, font=\scriptsize] at (axis cs:4.5,-1.5) {glass, $n = 1.5$};
    \node[black!60, font=\scriptsize] at (axis cs:7.7,-1.5) {vacuum (air), $n = 1$};
    % wavelength markers between crests
    \draw[<->, black!70] (axis cs:1.25,1.3) -- node[above, font=\scriptsize] {$\lambda_0$} (axis cs:2.25,1.3);
    \draw[<->, figR] (axis cs:3+1.25/1.5,1.3) -- node[above, font=\scriptsize] {$\lambda_0/n$} (axis cs:3+2.25/1.5,1.3);
    \draw[<->, black!70] (axis cs:6.75,1.3) -- node[above, font=\scriptsize] {$\lambda_0$} (axis cs:7.75,1.3);
  \end{axis}
\end{tikzpicture}
